% Paper 7 skeleton (DRAFT, workplace 2026-11-06; task V). Numbers come only from notes/paper7/out/tables.tex
% (generated by notes/paper7/tables.py from committed records). The text is finished at home.
\documentclass[10pt,a4paper]{article}
\usepackage[margin=2.2cm]{geometry}
\usepackage{amsmath,amssymb}
\usepackage{booktabs}
\usepackage[hidelinks]{hyperref}
\usepackage{url}
\setlength{\parskip}{4pt}
\setlength{\parindent}{0pt}
\sloppy

\title{Exact answers, no map, no mitigation:\\ user-written two-dimensional circuits through a placement-only OS on
IBM Heron r2 devices\\ \large DRAFT}
\author{TN Holdings\\ \small contact: \texttt{love.os.architect@proton.me}}
\date{2026 (draft of 6 November 2026, CET)}

\begin{document}
\maketitle

\begin{abstract}
% 1. Question: a user's plain-Qiskit 2D circuit with an exact answer, placed by translation without reading any
%    calibration-quality number, run without error mitigation: how close is the answer, and how is it shown?
% 2. What was run: practice problems (22 circuits x 2 devices, 2026-11-04), the daily record (from 2026-11-05),
%    the hard problems H01-H03 (from 2026-11-06). [numbers from tables (a)-(c)]
% 3. Declared verdicts first, retention beside them; the bias of an absolute criterion in both directions.
% 4. Negative: no correction works across problems (floor and mirror, declared; twirl, exploratory).
% 5. Not claimed: quantum advantage. Every answer is computed exactly on a classical machine (light cones of at most
%    33 qubits).
\end{abstract}

\section{Introduction: the no-map policy}
% Placement reads the coupling graph and failed / masked elements only; no quality numbers, no search.
% From paper 6 (ring self-check layer, 1D; doi:10.5281/zenodo.23883938) to the 2D user entrance (psf-os run2d).

\section{The OS's two-dimensional layer}
\subsection{Placement by translation}
\subsection{Floor, mirror and ring sentinel}
\subsection{The answer card}
\subsection{Writing $R_{ZZ}(\pm\pi/2)$ as one CZ}
% Before / after psf-os ac84182: CZ halved for P4 / P5 (results/practice-2026-11-04/README.md).

\section{Problems and answers}
\subsection{Problems after published experiments}
% P4: arXiv 2411.00765; P5: Kim et al., Nature 618, 500 (2023); P6: Qedma QESEM tutorial; daily pool; H01-H03.
\subsection{Sealing}
% Hashes committed before values: 12164be -> 50e412e.
\subsection{Two independent computations}

\section{Result 1: practice problems}
% % generated by notes/paper7/tables.py from committed records (DRAFT); do not edit by hand
\begin{table}[htbp]\centering\small
\caption{(a) Practice problems, 2026-11-04 (declared; raw, no correction). The declared criterion is an absolute difference; the same relative loss gives a larger difference when the exact answer is larger. R and the normalised error show the loss relative to the answer.}
\resizebox{\ifdim\width>\textwidth\textwidth\else\width\fi}{!}{\begin{tabular}{lllllllll}\toprule
device & problem & READS & AGREES & BEATS & M - M\_exact & R & norm. error & job \\ \midrule
ibm\_marrakesh & p4\_b0.485\_t02 & READS & AGREES & BEATS & -0.005 & 0.978 & 0.022 & dblibpd8mdtc73d66v90 \\
ibm\_marrakesh & p4\_b0.485\_t04 & READS & AGREES & BEATS & -0.006 & 0.972 & 0.028 & dblibpd8mdtc73d66v90 \\
ibm\_marrakesh & p4\_b0.485\_t08 & READS & AGREES & BEATS & -0.021 & 0.887 & 0.113 & dblibpd8mdtc73d66v90 \\
ibm\_marrakesh & p4\_b0.485\_t12 & READS & AGREES & BEATS & -0.022 & 0.862 & 0.138 & dblibpd8mdtc73d66v90 \\
ibm\_marrakesh & p4\_b0.485\_t16 & READS & AGREES & BEATS & -0.031 & 0.799 & 0.201 & dblibpd8mdtc73d66v90 \\
ibm\_marrakesh & p4\_b0.735\_t02 & READS & AGREES & BEATS & +0.024 & small answer & small answer & dblibpd8mdtc73d66v90 \\
ibm\_marrakesh & p4\_b0.735\_t04 & READS & AGREES & BEATS & +0.015 & small answer & small answer & dblibpd8mdtc73d66v90 \\
ibm\_marrakesh & p4\_b0.735\_t08 & READS & AGREES & BEATS & +0.013 & small answer & small answer & dblibpd8mdtc73d66v90 \\
ibm\_marrakesh & p4\_b0.735\_t12 & READS & AGREES & BEATS & +0.005 & small answer & small answer & dblibpd8mdtc73d66v90 \\
ibm\_marrakesh & p4\_b0.735\_t16 & READS & AGREES & BEATS & +0.001 & small answer & small answer & dblibpd8mdtc73d66v90 \\
ibm\_marrakesh & p5\_tx0.4\_t2 & READS & AGREES & DOES NOT BEAT & -0.046 & 0.950 & 0.050 & dblibpd8mdtc73d66v90 \\
ibm\_marrakesh & p5\_tx0.4\_t4 & READS & DIFFERS & DOES NOT BEAT & -0.063 & 0.933 & 0.067 & dblibpd8mdtc73d66v90 \\
ibm\_marrakesh & p5\_tx0.4\_t6 & READS & DIFFERS & DOES NOT BEAT & -0.060 & 0.933 & 0.067 & dblibpd8mdtc73d66v90 \\
ibm\_marrakesh & p5\_tx0.4\_t8 & READS & DIFFERS & BEATS & -0.070 & 0.922 & 0.078 & dblibpd8mdtc73d66v90 \\
ibm\_marrakesh & p5\_tx0.8\_t2 & READS & AGREES & BEATS & -0.024 & 0.963 & 0.037 & dblibpd8mdtc73d66v90 \\
ibm\_marrakesh & p5\_tx0.8\_t4 & READS & AGREES & BEATS & -0.030 & 0.941 & 0.059 & dblibpd8mdtc73d66v90 \\
ibm\_marrakesh & p5\_tx0.8\_t6 & READS & AGREES & BEATS & -0.021 & 0.951 & 0.049 & dblibpd8mdtc73d66v90 \\
ibm\_marrakesh & p5\_tx0.8\_t8 & READS & AGREES & BEATS & -0.048 & 0.893 & 0.107 & dblibpd8mdtc73d66v90 \\
ibm\_marrakesh & p5\_tx1.0\_t2 & READS & AGREES & BEATS & -0.010 & 0.977 & 0.023 & dblibpd8mdtc73d66v90 \\
ibm\_marrakesh & p5\_tx1.0\_t4 & READS & AGREES & BEATS & -0.000 & 0.998 & 0.002 & dblibpd8mdtc73d66v90 \\
ibm\_marrakesh & p5\_tx1.0\_t6 & READS & AGREES & BEATS & +0.009 & 1.054 & 0.054 & dblibpd8mdtc73d66v90 \\
ibm\_marrakesh & p5\_tx1.0\_t8 & READS & AGREES & BEATS & -0.004 & 0.976 & 0.024 & dblibpd8mdtc73d66v90 \\
ibm\_fez & p4\_b0.485\_t02 & READS & AGREES & BEATS & -0.018 & 0.917 & 0.083 & dblin89kbkrc73cnpcbg \\
ibm\_fez & p4\_b0.485\_t04 & READS & AGREES & BEATS & -0.027 & 0.882 & 0.118 & dblin89kbkrc73cnpcbg \\
ibm\_fez & p4\_b0.485\_t08 & READS & AGREES & BEATS & -0.039 & 0.785 & 0.215 & dblin89kbkrc73cnpcbg \\
ibm\_fez & p4\_b0.485\_t12 & READS & AGREES & BEATS & -0.041 & 0.740 & 0.260 & dblin89kbkrc73cnpcbg \\
ibm\_fez & p4\_b0.485\_t16 & READS & DIFFERS & BEATS & -0.060 & 0.612 & 0.388 & dblin89kbkrc73cnpcbg \\
ibm\_fez & p4\_b0.735\_t02 & READS & AGREES & BEATS & +0.014 & small answer & small answer & dblin89kbkrc73cnpcbg \\
ibm\_fez & p4\_b0.735\_t04 & READS & AGREES & BEATS & +0.010 & small answer & small answer & dblin89kbkrc73cnpcbg \\
ibm\_fez & p4\_b0.735\_t08 & READS & AGREES & BEATS & -0.001 & small answer & small answer & dblin89kbkrc73cnpcbg \\
ibm\_fez & p4\_b0.735\_t12 & READS & AGREES & BEATS & -0.013 & small answer & small answer & dblin89kbkrc73cnpcbg \\
ibm\_fez & p4\_b0.735\_t16 & READS & AGREES & BEATS & -0.018 & small answer & small answer & dblin89kbkrc73cnpcbg \\
ibm\_fez & p5\_tx0.4\_t2 & READS & AGREES & BEATS & -0.041 & 0.956 & 0.044 & dblin89kbkrc73cnpcbg \\
ibm\_fez & p5\_tx0.4\_t4 & READS & DIFFERS & DOES NOT BEAT & -0.061 & 0.936 & 0.064 & dblin89kbkrc73cnpcbg \\
ibm\_fez & p5\_tx0.4\_t6 & READS & DIFFERS & DOES NOT BEAT & -0.061 & 0.933 & 0.067 & dblin89kbkrc73cnpcbg \\
ibm\_fez & p5\_tx0.4\_t8 & READS & DIFFERS & BEATS & -0.056 & 0.938 & 0.062 & dblin89kbkrc73cnpcbg \\
ibm\_fez & p5\_tx0.8\_t2 & READS & AGREES & BEATS & -0.023 & 0.965 & 0.035 & dblin89kbkrc73cnpcbg \\
ibm\_fez & p5\_tx0.8\_t4 & READS & AGREES & BEATS & -0.021 & 0.958 & 0.042 & dblin89kbkrc73cnpcbg \\
ibm\_fez & p5\_tx0.8\_t6 & READS & AGREES & BEATS & -0.012 & 0.973 & 0.027 & dblin89kbkrc73cnpcbg \\
ibm\_fez & p5\_tx0.8\_t8 & READS & AGREES & BEATS & -0.017 & 0.961 & 0.039 & dblin89kbkrc73cnpcbg \\
ibm\_fez & p5\_tx1.0\_t2 & READS & AGREES & BEATS & -0.009 & 0.980 & 0.020 & dblin89kbkrc73cnpcbg \\
ibm\_fez & p5\_tx1.0\_t4 & READS & AGREES & BEATS & +0.006 & 1.025 & 0.025 & dblin89kbkrc73cnpcbg \\
ibm\_fez & p5\_tx1.0\_t6 & READS & AGREES & BEATS & +0.015 & 1.096 & 0.096 & dblin89kbkrc73cnpcbg \\
ibm\_fez & p5\_tx1.0\_t8 & READS & AGREES & BEATS & +0.012 & 1.072 & 0.072 & dblin89kbkrc73cnpcbg \\
\bottomrule\end{tabular}}
\end{table}

\begin{table}[htbp]\centering\small
\caption{(b) Daily record (declared; raw, no correction). The declared criterion is an absolute difference; the same relative loss gives a larger difference when the exact answer is larger. R and the normalised error show the loss relative to the answer.}
\resizebox{\ifdim\width>\textwidth\textwidth\else\width\fi}{!}{\begin{tabular}{llllllllll}\toprule
date & device & problem & AGREES & BEATS & M - M\_exact & R & norm. error & decay / CZ layer & job \\ \midrule
2026-11-05 & ibm\_fez & D02 & DIFFERS & BEATS & -0.147 & 0.816 & 0.184 & 0.0052 & dbm5hid8ft8s738u88g0 \\
2026-11-05 & ibm\_fez & FX-A & AGREES & BEATS & -0.043 & 0.732 & 0.268 & 0.0130 & dbm5hid8ft8s738u88g0 \\
2026-11-05 & ibm\_marrakesh & D02 & DIFFERS & BEATS & -0.111 & 0.860 & 0.140 & 0.0038 & dbm5hgd8mdtc73d7bang \\
2026-11-05 & ibm\_marrakesh & FX-A & AGREES & BEATS & -0.024 & 0.846 & 0.154 & 0.0070 & dbm5hgd8mdtc73d7bang \\
2026-11-06 & ibm\_marrakesh & D03 & REFUSED & - & - & - & - & - & not run (no copy qualifies) \\
2026-11-06 & ibm\_marrakesh & FX-A & AGREES & BEATS & -0.008 & 0.949 & 0.051 & 0.0017 & dbmqhbu74qds73dsrde0 \\
2026-11-07 & ibm\_fez & D04 & AGREES & BEATS & +0.010 & 1.022 & 0.022 & -0.0005 & dbn5ludt7fqc73aivq60 \\
2026-11-07 & ibm\_fez & FX-A & AGREES & BEATS & -0.038 & 0.763 & 0.237 & 0.0121 & dbn5ludt7fqc73aivq60 \\
\bottomrule\end{tabular}}
\end{table}

\begin{table}[htbp]\centering\small
\caption{(c) FX-A day to day (the same circuit as practice p4\_b0.485\_t12). The declared criterion is an absolute difference; the same relative loss gives a larger difference when the exact answer is larger. R and the normalised error show the loss relative to the answer.}
\resizebox{\ifdim\width>\textwidth\textwidth\else\width\fi}{!}{\begin{tabular}{llllll}\toprule
date & device & M - M\_exact & R & AGREES & job \\ \midrule
2026-11-04 (practice) & ibm\_marrakesh & -0.022 & 0.862 & AGREES & dblibpd8mdtc73d66v90 \\
2026-11-04 (practice) & ibm\_fez & -0.041 & 0.740 & AGREES & dblin89kbkrc73cnpcbg \\
2026-11-05 & ibm\_fez & -0.043 & 0.732 & AGREES & dbm5hid8ft8s738u88g0 \\
2026-11-05 & ibm\_marrakesh & -0.024 & 0.846 & AGREES & dbm5hgd8mdtc73d7bang \\
2026-11-06 & ibm\_marrakesh & -0.008 & 0.949 & AGREES & dbmqhbu74qds73dsrde0 \\
2026-11-07 & ibm\_fez & -0.038 & 0.763 & AGREES & dbn5ludt7fqc73aivq60 \\
\bottomrule\end{tabular}}
\end{table}

\begin{table}[htbp]\centering\small
\caption{(d) Corrections compared (declared 11-04; exploratory 11-05 marked so: corrected MSE over the qubits kept by rule R, raw over all; AGREES after computed here with the same 0.05). The declared criterion is an absolute difference; the same relative loss gives a larger difference when the exact answer is larger. R and the normalised error show the loss relative to the answer.}
\resizebox{\ifdim\width>\textwidth\textwidth\else\width\fi}{!}{\begin{tabular}{lllllllll}\toprule
kind & device & problem & raw MSE & corrected MSE & ratio & verdict & AGREES after & job \\ \midrule
declared (11-04) & ibm\_fez & P6 theta\_x 0.53 & 0.0419 & 0.0034 & 0.082 & HELPS & AGREES & dblgrsm74qds73dqfbqg \\
declared (11-04) & ibm\_fez & P6 theta\_x 0.3 & 0.0439 & 0.0008 & 0.018 & HELPS & AGREES & dblgrsm74qds73dqfbqg \\
declared (11-04) & ibm\_fez & P6 theta\_x 0.8 & 0.0205 & 0.0083 & 0.403 & HELPS & DIFFERS & dblgrsm74qds73dqfbqg \\
declared (11-04) & ibm\_marrakesh & P6 theta\_x 0.53 & 0.0897 & 0.0169 & 0.189 & HELPS & DIFFERS & dblgv8pkbkrc73cnn6qg \\
declared (11-04) & ibm\_marrakesh & P6 theta\_x 0.3 & 0.0875 & 0.0092 & 0.105 & HELPS & DIFFERS & dblgv8pkbkrc73cnn6qg \\
declared (11-04) & ibm\_marrakesh & P6 theta\_x 0.8 & 0.0450 & 0.0072 & 0.161 & HELPS & AGREES & dblgv8pkbkrc73cnn6qg \\
exploratory, not declared (11-05, run) & ibm\_fez & D02 corr\_raw & 0.0343 & 0.0117 & 0.341 & not judged & AGREES & dbm6471kbkrc73cou7e0 \\
exploratory, not declared (11-05, run) & ibm\_fez & D02 corr\_tw & 0.0343 & 0.0316 & 0.921 & not judged & AGREES & dbm6471kbkrc73cou7e0 \\
exploratory, not declared (11-05, run) & ibm\_fez & FX-A corr\_raw & 0.0027 & 0.0085 & 3.104 & not judged & DIFFERS & dbm6471kbkrc73cou7e0 \\
exploratory, not declared (11-05, run) & ibm\_fez & FX-A corr\_tw & 0.0027 & 0.0073 & 2.665 & not judged & AGREES & dbm6471kbkrc73cou7e0 \\
exploratory, not declared (11-05, run2) & ibm\_fez & D02 corr\_raw & 0.0372 & 0.0133 & 0.357 & not judged & AGREES & dbm6jfd8mdtc73d7dit0 \\
exploratory, not declared (11-05, run2) & ibm\_fez & D02 corr\_tw & 0.0372 & 0.0346 & 0.929 & not judged & AGREES & dbm6jfd8mdtc73d7dit0 \\
exploratory, not declared (11-05, run2) & ibm\_fez & D02 corr\_btw & 0.0372 & 0.0296 & 0.796 & not judged & AGREES & dbm6jfd8mdtc73d7dit0 \\
exploratory, not declared (11-05, run2) & ibm\_fez & FX-A corr\_raw & 0.0048 & 0.0325 & 6.755 & not judged & DIFFERS & dbm6jfd8mdtc73d7dit0 \\
exploratory, not declared (11-05, run2) & ibm\_fez & FX-A corr\_tw & 0.0048 & 0.0222 & 4.610 & not judged & DIFFERS & dbm6jfd8mdtc73d7dit0 \\
exploratory, not declared (11-05, run2) & ibm\_fez & FX-A corr\_btw & 0.0048 & 0.0198 & 4.109 & not judged & DIFFERS & dbm6jfd8mdtc73d7dit0 \\
exploratory, not declared (11-05, run) & ibm\_marrakesh & D02 corr\_raw & 0.0125 & 0.0043 & 0.343 & not judged & AGREES & dbm644hkbkrc73cou76g \\
exploratory, not declared (11-05, run) & ibm\_marrakesh & D02 corr\_tw & 0.0125 & 0.0044 & 0.357 & not judged & AGREES & dbm644hkbkrc73cou76g \\
exploratory, not declared (11-05, run) & ibm\_marrakesh & FX-A corr\_raw & 0.0030 & 0.0168 & 5.609 & not judged & DIFFERS & dbm644hkbkrc73cou76g \\
exploratory, not declared (11-05, run) & ibm\_marrakesh & FX-A corr\_tw & 0.0030 & 0.0065 & 2.164 & not judged & DIFFERS & dbm644hkbkrc73cou76g \\
exploratory, not declared (11-05, run2) & ibm\_marrakesh & D02 corr\_raw & 0.0155 & 0.0032 & 0.207 & not judged & AGREES & dbm6jcl8mdtc73d7dio0 \\
exploratory, not declared (11-05, run2) & ibm\_marrakesh & D02 corr\_tw & 0.0155 & 0.0053 & 0.346 & not judged & AGREES & dbm6jcl8mdtc73d7dio0 \\
exploratory, not declared (11-05, run2) & ibm\_marrakesh & D02 corr\_btw & 0.0155 & 0.0055 & 0.359 & not judged & AGREES & dbm6jcl8mdtc73d7dio0 \\
exploratory, not declared (11-05, run2) & ibm\_marrakesh & FX-A corr\_raw & 0.0021 & 0.0268 & 12.499 & not judged & DIFFERS & dbm6jcl8mdtc73d7dio0 \\
exploratory, not declared (11-05, run2) & ibm\_marrakesh & FX-A corr\_tw & 0.0021 & 0.0091 & 4.249 & not judged & DIFFERS & dbm6jcl8mdtc73d7dio0 \\
exploratory, not declared (11-05, run2) & ibm\_marrakesh & FX-A corr\_btw & 0.0021 & 0.0105 & 4.870 & not judged & DIFFERS & dbm6jcl8mdtc73d7dio0 \\
\bottomrule\end{tabular}}
\end{table}
 -- table (a)

\section{Result 2: the daily record}
% tables (b), (c); every night and device a row, failures included.

\section{Verdict and retention}
% notes/reporting/verdict_and_retention.md; footnote on the absolute criterion; decay per CZ layer by family.

\section{A negative result: correction}
% table (d): declared floor-and-mirror correction (2026-11-04); exploratory twirl (2026-11-05, not declared);
% after the fact: the loss next to dropped qubits (marked "after the fact").

\section{Reproduction}
% reproduce.py (9 of 9 the same), results/ENVIRONMENT.md, three levels.

\section{Limits}
% Devices, days, hours; P6's outside comparison (our unmitigated gaps larger than the tutorial's on 2026-11-04);
% patches placeable on one device only.

\section{Conclusion}

\section*{Acknowledgements}
We acknowledge the use of IBM Quantum services for this work. The views expressed are those of the authors and do not
reflect the official policy or position of IBM or the IBM Quantum team.

\section*{Data and code}
% TOQB and psf-os repositories (commit IDs), raw counts and job IDs in results/; reproduce.py.

% \bibliographystyle{unsrt}
% \bibliography{paper7}
\end{document}
